\section{Background}

\subsection{A relative valence--arousal coordinate}

The circumplex model represents affect along two continuous dimensions:
\emph{valence}, from unpleasant to pleasant, and \emph{arousal}, from low to
high activation~\cite{russell1980,posner2005}. Thayer's account of energetic and
tense arousal offers a complementary vocabulary for changes in activation and
mood~\cite{thayer1991}. These models are useful here as a compact interaction
space, not as a complete account of emotion. Both dimensional and discrete
models appear in music-emotion research, with results sensitive to model and
annotations~\cite{eerola2013}.

Figure~\ref{fig:affect-space} is an independently drawn conceptual synthesis:
it places the continuous and categorical affect vocabularies beside the
strictly group-level HRV-frequency context that motivates an exploratory
diagnostic signal.

\begin{figure*}[t]
  \centering
  \includegraphics[width=\textwidth]{fig1_affect_space.pdf}
  \caption{\textbf{Conceptual foundations.} Panels (a)--(b) are adapted and
  independently redrawn from Helmholz et al.~\cite{helmholz2019}, Figure~1,
  which operationalizes Russell's circumplex~\cite{russell1980} and Thayer's
  mood--arousal model~\cite{thayer1991}; no source pixels are reused. Panel (c)
  is an original, data-free schematic of the group-level frequency distinction
  discussed by Balaji et al.~\cite{balaji2025}; no curves, observations, or
  source artwork are reproduced or adapted. Illustrations only---not Vibe
  measurements or classifiers.}
  \Description{Three square conceptual panels show a continuous
  valence-arousal plane, a categorical affect vocabulary, and two synthetic
  coherence-frequency profiles explicitly marked as schematic rather than
  data.}
  \label{fig:affect-space}
\end{figure*}

Vibe Shuffle borrows only the two-axis coordinate structure. It does not output
the emotion words shown in Figure~\ref{fig:affect-space} or treat them as
ground-truth labels. The coherence profiles in panel (c) are likewise not the
prototype's diagnostic signal.

Vibe Shuffle places the listener at a point $(V,A)$ in this space and divides
the plane into four quadrants. Crucially, the point describes change relative
to that listener's own reference measurements. The center is the personal
baseline; evidence of more positive facial display moves the point right, while
evidence of greater cardiac activation moves it upward. The system therefore
asks a deliberately narrow question---which direction did the observed signals
move?---rather than attempting to name or recover the listener's full emotional
state. This distinction matters because facial movements are not unique
fingerprints of specific emotions and vary with person and
context~\cite{barrett2019}. Facial judgments can also shift with body context,
and Valence--Arousal relations vary across people and
cultures~\cite{aviezer2012,kuppens2017}.

\subsection{Connecting listener state to music}

Music can be placed in a parallel two-dimensional space. Song \emph{Valence}
describes conveyed positivity, whereas \emph{Energy} describes perceived
intensity and activity. Prior emotion-oriented music interfaces have used these
attributes as musical counterparts to affective valence and
arousal~\cite{helmholz2019}. Vibe Shuffle adopts that practical bridge: listener
Valence is compared with song Valence, and listener Arousal is compared with
song Energy. The alignment is an engineering assumption, not a claim that
musical Energy and physiological Arousal are the same construct.
Music-emotion recognition also uses continuous coordinates, but its annotations
remain subjective~\cite{yang2012}.

The shared plane makes the selection rule easy to inspect. High and low values
on each axis produce four broad regions, and the listener's relative position
determines the target region. Within that region, distance in the continuous
plane distinguishes closer from more distant songs. The quadrants simplify
communication; the underlying coordinate preserves degree and direction.

\subsection{Heart-rate variability as relative evidence}

Heart-rate variability (HRV) describes variation between successive heart
beats. The root mean square of successive differences (RMSSD) is a widely used
time-domain summary of short-term RR-interval variation~\cite{shaffer2017}.
Heart rate and RMSSD depend on respiration, movement, posture, fitness,
medication, sensor quality, and other factors; neither measure identifies a
particular emotion on its own. Good-practice guidance therefore emphasizes
standardized conditions, artifact review, matched durations, and respiration
assessment~\cite{laborde2017}.

Vibe Shuffle consequently uses cardiac measurements only as
baseline-relative evidence for the arousal axis. A heart rate above the
listener's reference and an RMSSD below that reference contribute upward
evidence; changes in the opposite direction lower it. The mapping is transparent
and testable, but remains a prototype heuristic rather
than a clinical or psychological measurement.
